\chapter{いつ探索し、いつ決めるか：\nt{NORA}を加える}
\chaptermark{\nt{NORA}を加える}
\label{sec:nora}
\begin{chaptergoals}
\item 一つのゲインのつまみがニューロンをアナログ増幅器からコンパレータに変える仕組み；
\item なぜゲインを上げて役立つのは増幅器の後で生じるノイズに対してだけなのか；
\item 重みを変えずにゲインが選択ネットワークを「思案中」から「決定済み」にする仕組み；
\item なぜゲインは逆温度として働き、報酬はゲインに対して逆U字になるのか。
\end{chaptergoals}

\section{足りないもの}

第\ref{sec:dofa}章でネットワークはノイズのおかげで学習した。活動のランダムなずれが探索であり、\nt{DOPA}はそれでよくなったかを伝えた。しかしそこでは一つのパラメータが未設定のまま残った。\emph{どれだけ}探索するかである。ノイズが強すぎると、ネットワークはあちこちに振れ、すでに知っていることを使わない。弱すぎると、毎回同じものを選び、近くのどこかがよくなったことに気づかない。強化学習ではこれを活用と探索の選択と呼び、どのアルゴリズムにもそのためのつまみがある。第\ref{sec:neurontypes}章で、脳ではこのつまみを回すのは\nt{NORA}だと仮定した。

ヒトのノルアドレナリンニューロンは数万個で（表~\ref{tab:types}）、ほぼすべてが一つの小さな集団に集まっている。その出力は脳のほぼ全体に広がるが、集団の部分ごとに、ある程度異なる領域を担当していることが分かっている~\cite{Poe2020}。これは\nt{DOPA}よりさらに細いブロードキャストチャネルである。動作モードは二つある~\cite{AstonJones2005}。\emph{持続}モードでは、ニューロンは毎秒1回未満から数回の発火率で規則的にスパイクを出し、この発火率は覚醒度とともにゆっくり変わる。\emph{相動}モードでは、現在の課題にとって重要な出来事に、ほぼ意思決定の瞬間に短いバーストで応答する。便利だが不正確なテストポイントが瞳孔径である。瞳孔径はこれらのニューロンの活動とともに変わるが、他の領域の活動~\cite{Joshi2016}や皮質のコリン作動性出力~\cite{Reimer2016}とともにも変わる。瞳孔が示すのは全体的な興奮レベルであり、\nt{NORA}だけではない。

\section{ゲイン}

測定されているのは次のことである。ノルアドレナリンはニューロンの背景活動を刺激への応答より強く抑え、信号対背景比が上がる~\cite{Foote1975NE}。個々のニューロンの特性の傾きが文字どおり係数倍されることは、実験からは導かれない。第\ref{sec:neurontypes}章と同じく、この効果を再現する簡単なモデルをとる。ノルアドレナリンは受け手の伝達特性の傾きを変える~\cite{ServanSchreiber1990}。すると弱い閾値下の入力（背景）はさらに小さな出力を、強い入力はさらに大きな出力を生む。ニューロンは発火率で応答するとする。
\begin{empheq}[box=\keybox]{equation}
r \;=\; F\bigl(g\,(u-\theta)\bigr), \qquad F(z)=\frac{r_{\max}}{1+e^{-z}} ,
\label{eq:gain}
\end{empheq}
ここで$u$は入力電流、$\theta$は閾値、$g$はゲインで、モデルではこれを\nt{NORA}が制御する。$g$が小さいと、発火率は広い範囲で入力になめらかに従う。ニューロンは\emph{アナログ増幅器}である。$g$が大きいと、閾値より上のほぼどんな入力も最大発火率を、下なら0を生む。ニューロンは\emph{コンパレータ}、ほぼディジタル素子である（図~\ref{fig:gaincurves}）。一つのつまみが、電子回路なら別々の回路が担う二つのモードの間で同じニューロンを切り替える。

\begin{figure}[t]
\centering
\begin{tikzpicture}[>=Stealth,x=1.1cm,y=2.4cm]
\draw[->,gray!70!black] (-3.3,0) -- (3.4,0) node[below left,black] {\small $u$};
\draw[->,gray!70!black] (-3.3,0) -- (-3.3,1.2) node[left,black] {\small $r$};
\draw[densely dashed,gray!60!black] (0,0) -- (0,1.1);
\node[below,font=\small] at (0,0) {$\theta$};
\draw[densely dotted,gray!60!black] (-3.3,1) -- (3.2,1) node[right,black,font=\small] {$r_{\max}$};
\draw[very thick,blue!60!black] plot[domain=-3.2:3.2,samples=60] (\x,{1/(1+exp(-0.8*\x))});
\draw[very thick,orange!85!black] plot[domain=-3.2:3.2,samples=60] (\x,{1/(1+exp(-2*\x))});
\draw[very thick,red!70!black] plot[domain=-3.2:3.2,samples=120] (\x,{1/(1+exp(min(-8*\x,9)))});
\node[blue!60!black,font=\small,anchor=west] at (0.5,0.12) {小さい$g$：増幅器};
\node[red!70!black,font=\small,anchor=east] at (-0.3,0.85) {大きい$g$：コンパレータ};
\end{tikzpicture}
\caption{ゲイン$g$の三つの値での伝達特性~\eqref{eq:gain}。}
\label{fig:gaincurves}
\end{figure}

大きなゲインはノイズに対して役立つか。いつもではない。二つの入力が$\Delta u$だけ異なり、どちらが大きいかを言わなければならないとする。ノイズ$\sigma_{\text{pre}}$が増幅器の\emph{前で}入力に加わるなら、信号もノイズも増幅され、その比は変わらない。ノイズ$\sigma_{\text{post}}$が増幅器の\emph{後で}加わるなら（第\ref{sec:code}章のように、信頼性の低いシナプスやネットワーク自体のランダムなスパイクから）、SN比
\begin{empheq}[box=\keybox]{equation}
d' \;=\; \frac{g\,\Delta u}{\sqrt{g^2\sigma_{\text{pre}}^2+\sigma_{\text{post}}^2}}
\label{eq:gainsnr}
\end{empheq}
は、$g\sigma_{\text{pre}}$が$\sigma_{\text{post}}$に並ぶまで$g$とともに増える~\cite{ServanSchreiber1990}。

\begin{keyblock}
\begin{proposition}[ゲインが役立つとき]
\label{prop:gainnoise}
ゲインを上げることで\nt{NORA}が入力の識別を改善できるのは、増幅器の\emph{後で}、ネットワーク自体の中で生じるノイズに対してだけである。入力と一緒に来たノイズはゲインでは除けない。
\end{proposition}
\end{keyblock}

\section{二者択一でのゲイン}

第\ref{sec:gaba}章の選択に戻ろう。入力$I_1$、$I_2$をもつ\nt{GLUT}ニューロンの二つの群が、強さ$\beta$で互いに抑制し合う。今度は各群にゲイン$g$がある。選択の方程式~\eqref{eq:wta}の括弧内の式が$g$倍される。両方の群が動いている間、定常発火率は$r_1=g(I_1-\beta r_2)$、$r_2=g(I_2-\beta r_1)$から求まる。これらを足し引きすると、和$r_1+r_2=g(I_1+I_2)/(1+g\beta)$と差$r_1-r_2=g(I_1-I_2)/(1-g\beta)$が得られる。その比は、ネットワークが一方の入力をどれだけはっきり好むかを表す。
\begin{empheq}[box=\keybox]{equation}
\frac{r_1-r_2}{r_1+r_2} \;=\; \varepsilon\,\frac{1+g\beta}{1-g\beta}, \qquad \varepsilon=\frac{I_1-I_2}{I_1+I_2}, \quad g\beta<1 .
\label{eq:wtagain}
\end{empheq}
入力の小さな差$\varepsilon$は$(1+g\beta)/(1-g\beta)$倍に増幅され、この係数は$g\beta$が1に近づくと際限なく大きくなる。$g\beta>1$では、命題~\ref{prop:wta}と同じく、両方の群が動く状態は不安定である。一方が勝ち、他方は沈黙する（図~\ref{fig:pitchfork}）。

\begin{figure}[t]
\centering
\begin{tikzpicture}[>=Stealth,x=3.2cm,y=1.5cm]
\draw[->,gray!70!black] (0,0) -- (2.15,0) node[below left,black] {\small $g\beta$};
\draw[->,gray!70!black] (0,-1.25) -- (0,1.35) node[above,black] {\small $(r_1-r_2)/(r_1+r_2)$};
\draw[gray!60!black] (1,-0.04) -- (1,0.04); \node[below right,font=\small] at (1,0) {$1$};
\node[left,font=\scriptsize] at (0,1) {$1$}; \node[left,font=\scriptsize] at (0,-1) {$-1$};
\draw[very thick,blue!60!black] plot[domain=0:0.905,samples=60] (\x,{0.05*(1+\x)/(1-\x)}) -- (1,1) -- (2.05,1);
\draw[very thick,blue!60!black] (1.105,-1) -- (2.05,-1);
\draw[thick,densely dashed,gray!60!black] (1,0) -- (2.05,0);
\node[font=\small,align=center] at (0.45,-0.62) {思案中：\\両方の群が\\動いている};
\node[font=\small,align=center] at (1.55,0.45) {決定済み：\\一方だけが動く};
\node[font=\small,blue!60!black,anchor=south west] at (1.05,-0.97) {弱い入力が先に勝つと、そのまま保たれる};
\end{tikzpicture}
\caption{ゲインを変えたときの二者択一。入力の差は$\varepsilon=0.05$。$g\beta>1$では両方の決定が安定で（実線。弱い入力の勝利は$g\beta>(1+\varepsilon)/(1-\varepsilon)\approx1.1$で）、ネットワークはすでに下した決定を保つ。対称な状態（破線）は不安定である。}
\label{fig:pitchfork}
\end{figure}

\begin{keyblock}
\begin{proposition}[ゲインが選択のモードを切り替える]
\label{prop:gainwta}
相互抑制$\beta$をもつ選択ネットワークで、ゲイン$g$はネットワークを境界$g\beta=1$の向こうへ移す。その下ではネットワークは「思案中」である。両方の群が動き、入力の差は\eqref{eq:wtagain}に従って増幅される。その上ではネットワークは「決定済み」である。一方の群だけが動き、決定は保たれる。$g$を変えることで、\nt{NORA}は重みに一つも触れずにネットワークをこの二つのモードの間で切り替えられる。
\end{proposition}
\end{keyblock}

\section{温度としてのゲイン}

次に、ネットワーク自体の中、増幅器の後で生じるノイズを選択に加えよう。どちらの群が勝つかは、量$g(u_A-u_B)+\eta$の符号で決まる。ここで$u_A-u_B$は入力の差、つまり第\ref{sec:dofa}章で学習した重みの差であり、$\eta$はスケール$\sigma$のノイズである。ノイズがロジスティック分布に従うなら（ガウス分布でも曲線はほぼ同じ）、$A$を選ぶ確率は
\begin{empheq}[box=\keybox]{equation}
P(A) \;=\; \frac{1}{1+e^{-g\,(u_A-u_B)/\sigma}} .
\label{eq:softmax}
\end{empheq}
プログラマーならここに温度$T=\sigma/g$のソフトマックス選択を見てとるだろう。ゲインは逆温度である。$g$が小さいと、明らかによい選択肢でも悪いほうよりわずかに多く選ばれるだけである。ネットワークは多く探索する。$g$が大きいと、既知の最良の選択肢がほぼ常に選ばれる。ネットワークは知っていることを活用する。

\eqref{eq:wtagain}と\eqref{eq:softmax}の$g$が何かを明確にしておく。これは選択の瞬間における、現在の課題の入力の差に対する\emph{実効}ゲインである。\nt{NORA}の二つのモードはこれに逆向きに作用する。決定の瞬間の相動性バーストはこれを上げ、ネットワークは選択を固める。逆に高い持続性活動は探索と一緒に現れる。ヒトでは当て推量の選択の前にベースラインの瞳孔が広く~\cite{JepmaNieuwenhuis2011}、ラットでは前帯状皮質への\nt{NORA}入力が選択をランダムなモードに移す~\cite{Tervo2014stochastic}。したがって高い持続性\nt{NORA}には\emph{小さな}実効$g$が、バーストには大きな実効$g$が対応する。

\section{どれだけ探索するか}

どのゲインが最もよいか。モデルで確かめた。ネットワークは\eqref{eq:softmax}に従って二つの行動の一方を選ぶ。どちらもある確率で報酬をもたらし、ネットワークは第\ref{sec:dofa}章のように、選んだ行動の報酬からその確率を推定する。ときどき世界が変わり、二つの確率が引き直される。変わるのは選んだ行動の場合もあれば、ネットワークが長く試していない行動の場合もある。後者についてネットワークが知る方法は探索しかない。図~\ref{fig:invertedU}は、一定の$g$をいろいろ変えたとき、ネットワークが最良の可能な報酬のどれだけの割合を得るかを示す。

\begin{figure}[t]
\centering
\begin{tikzpicture}[>=Stealth,x=1.2cm,y=1cm]
\draw[->,gray!70!black] (0,0) -- (7.6,0) node[below left,black] {\small $g$};
\draw[->,gray!70!black] (0,0) -- (0,4.4) node[above,black,font=\small] {最良の報酬に対する割合};
\foreach \x/\l in {0/0.5,1/1,2/2,3/4,4/8,5/16,6/32,7/64}{\draw[gray!60!black] (\x,-0.06) -- (\x,0.06); \node[below,font=\scriptsize] at (\x,0) {\l};}
\foreach \v/\y in {0.8/0.8,0.9/2.4,1.0/4.0}{\draw[gray!60!black] (-0.06,\y) -- (0.06,\y); \node[left,font=\scriptsize] at (0,\y) {\v};}
\draw[very thick,blue!60!black] plot[mark=*,mark size=1.3pt] coordinates {(0,0.368) (1,0.880) (2,1.728) (3,2.784) (4,3.456) (5,3.616) (6,3.488) (7,3.264)};
\draw[very thick,red!70!black] plot[mark=*,mark size=1.3pt] coordinates {(0,0.368) (1,0.672) (2,1.184) (3,1.840) (4,2.176) (5,1.920) (6,1.456) (7,1.104)};
\node[blue!60!black,font=\small,anchor=south] at (5,3.7) {穏やかな世界};
\node[red!70!black,font=\small,anchor=north] at (5.1,1.25) {速く変わる世界};
\draw[->,thick,blue!60!black] (5,3.25) -- (5,3.55);
\draw[->,thick,red!70!black] (4,1.8) -- (4,2.1);
\node[font=\small\itshape,gray!35!black,anchor=west] at (0.15,2.3) {やみくもに探す};
\node[font=\small\itshape,gray!35!black] at (6.45,2.55) {変化に気づかない};
\end{tikzpicture}
\caption{一定のゲイン$g$での、最良の可能な報酬に対する割合（$g$の目盛りは対数）。青の曲線は世界が平均1000試行に1回変わる場合、赤は20試行に1回の場合。矢印は最良の$g$を示す。速く変わる世界ではそれは半分になる。4000試行の実行60回の平均。}
\label{fig:invertedU}
\end{figure}

$g=0.5$ではネットワークは知識をほとんど使わず、可能な報酬のおよそ4分の3を得る。$g$が非常に大きいと変化を見逃す。最良の値は、世界が1000試行に1回変わるとき$g=16$（割合$0.976$）、20試行に1回変わるとき$g=8$（$0.886$）である。

\begin{keyblock}
\begin{proposition}[逆U字]
\label{prop:explore}
ネットワークが集める報酬は、ゲインに対して逆U字形に依存する。$g$が小さすぎると試行を探索に費やし、大きすぎると変化に気づかない。最良の$g$は、世界が速く変わるほど小さい。
\end{proposition}
\end{keyblock}

逆U字は実験でも観察される。動物が課題を最もよく解くのは、\nt{NORA}ニューロンの持続性活動が中程度で、相動性応答が明瞭なときであり、持続性活動が高いと気が散り、課題の間を切り替える~\cite{AstonJones2005}。これは前節のモードの違いと整合する。決定の瞬間の相動性バーストは、選ぶべきときにちょうど大きな実効$g$を与える。バーストのない高い持続性活動は、無関係な入力も含めて何でも増幅するので、現在の課題にとっての実効$g$は下がる。これが探索である。

\section{誰がつまみを回すのか}

最良のゲインが世界の変化の速さに依存するなら、それを調整できるとよい。しかし\nt{NORA}ニューロンは世界が変わったことをどうして知るのか。自然な手がかりは\emph{驚き}である。予測誤差がふだんより大きくなった。二種類の不確実性を区別することが大切である。報酬が単にランダムなら、誤差は常に大きく、それは予期されている。一方、誤差がふだんの大きさに比べて急に増えたなら、何かが変わったのである。ある仮説によれば、\nt{NORA}が伝えるのはまさにこの予期しない不確実性であり、予期された不確実性は\nt{ACET}が伝える~\cite{YuDayan2005}。

この信号で何をするか。ゲインを下げて探索を増やしてもよい。古くなった推定を速く忘れるよう学習率を上げてもよい。実験では、急な変化の後に人はより速く学び、そのとき瞳孔が広がる~\cite{Nassar2012}。瞳孔は\nt{NORA}だけでなく\nt{ACET}の指標でもあることを思い出そう~\cite{Reimer2016}。両方を同時にしてもよい。さらに別の仮説によれば、\nt{NORA}の相動性バーストは\emph{割り込み}として働く。プロセッサの共通の割り込み線のように、ネットワークが現在の状態をリセットし、新しい状況に合わせて組み直す合図である~\cite{BouretSara2005}。

見つからなかったことも正直に述べておく。モデルで二つの簡単な調整法を試した。平滑化した誤差が長期平均を超えたら$g$を下げる方法と、学習率を上げる方法である。長く静止したり頻繁に変わったりする世界で、どちらの方法も最良の設定では最良の報酬の$0.926$--$0.927$を得た。これは最良の一定の$g$（$g=8$で$0.925$）や、一定だが十分大きな学習率（$0.928$）とほとんど同じであり、設定が悪いとそれより少なかった。図~\ref{fig:invertedU}の曲線は頂上付近でなだらかで、このような世界では調整することがほとんどない。調整が割に合うのは、モードごとに大きく異なる設定が要る場合のはずである。\nt{NORA}がどんな制御則を実装しているのかは、まだ確立していない。

\section{まとめ}

\begin{table}[t]
\centering
\footnotesize
\setlength{\tabcolsep}{4pt}
\renewcommand{\arraystretch}{1.15}
\begin{tabular}{>{\raggedright}p{3.0cm}>{\raggedright}p{4.4cm}>{\raggedright\arraybackslash}p{6.0cm}}
\toprule
\nt{NORA}がすること & 方法 & 分かっていること \\
\midrule
ニューロンの応答の傾きを変える & \eqref{eq:gain}のゲイン$g$ & SN比の向上は測定済み。乗算モデルは単純化 \\
選択を切り替える & ネットワークを$g\beta=1$の向こうへ移す & モデル~\eqref{eq:wtagain}から導かれる。閾値の直接測定はない \\
どれだけ探索するかを決める & $g$は逆温度~\eqref{eq:softmax}。持続性は小さな実効$g$（探索）、バーストは大きな実効$g$（決定） & 逆U字と二つの動作モードは測定済み。探索との関係は十分な根拠のある仮説 \\
変化を伝える & 予期しない不確実性 & 変化の後に瞳孔と学習率が増すことは測定済み。それが\nt{NORA}の信号であることは仮説 \\
割り込み、リセットする & 予期しない出来事への相動性バースト & 重要な出来事へのバーストは測定済み。「割り込み」は仮説 \\
\bottomrule
\end{tabular}
\caption{\nt{GLUT}、\nt{GABA}、\nt{DOPA}のネットワークに\nt{NORA}が加えるもの。}
\label{tab:nora}
\end{table}

\nt{DOPA}はネットワークに重みを\emph{どこへ}変えるかを伝える。\nt{NORA}は重みを一つも変えないが、ネットワークがすでに知っていることを\emph{どう}使うかを決める（表~\ref{tab:nora}）。一つのつまみ、ゲインが、ニューロンを増幅器からコンパレータへ、選択ネットワークを「思案中」から「決定済み」へ、選択を探索から活用へと変える。\nt{NORA}が世界の変化の速さをどう知るのかは未解決の問題である。

\section{問題}

\begin{exercise}[\lvlA]
\label{ex:08:1}
\emph{脳全体に一つのつまみ。}(a)~表~\ref{tab:types}から、\nt{NORA}ニューロン1個あたり脳のニューロンが何個あるかを見積もれ。(b)~増幅器の前のノイズが$\sigma_{\text{pre}}=0.5$、後のノイズが$\sigma_{\text{post}}=4$である。\eqref{eq:gainsnr}の$d'$が上限の$90\,\%$に達するのはどの$g$か。
\end{exercise}

\begin{exercise}[\lvlB]
\label{ex:08:2}
\emph{ゲイン付きの選択。}$\tau\,\dd r_1/\dd t=-r_1+g[I_1-\beta r_2]_+$、$\tau\,\dd r_2/\dd t=-r_2+g[I_2-\beta r_1]_+$、$\tau=20\ms$。(a)~$g\beta=0.5$と$0.9$で、$r_1-r_2$が減衰する時間はいくらか。(b)~$\varepsilon=0.05$のとき、それより下で両方の群が動き、それより上で弱い入力の勝利が安定になる$g\beta$を求めよ。
\end{exercise}

\begin{exercise}[\lvlA]
\label{ex:08:3}
\emph{ノイズからソフトマックスへ。}$K$個の群はそれぞれ自分のガンベルノイズ$P(\eta_k<z)=\exp(-e^{-z/\sigma})$を受け、$gu_k+\eta_k$が最大のものが勝つ。$P(k)$を求めよ。$u_A-u_B=0.1$、$\sigma=1$の二つの選択肢のうちよいほうが$90\,\%$の確率で選ばれるのはどの$g$か。
\end{exercise}

\begin{exercise}[\lvlB]
\label{ex:08:4}
\emph{最良のゲインの見積もり。}図~\ref{fig:invertedU}のモデルで、変化後の報酬確率は$[0,1]$上で一様である。ネットワークは悪いほうの行動を$e^{g\Delta}$試行に1回試す。これを$1/h$に等しいとおくと$g^*\approx\ln(1/h)/\bar\Delta$となる。$\bar\Delta=\langle|p_A-p_B|\rangle$と、$h=0.001$、$0.01$、$0.05$での$g^*$を求めよ。
\end{exercise}

\begin{exercise}[\lvlB]
\label{ex:08:5}
\emph{シミュレーション：ゲインと変化の速さ。}\texttt{models/\allowbreak nora\_explore.py}のコピー（ファイルを現在のフォルダに書き出す）で、$h$が$0.001$から$0.2$までの$g^*(h)$を求めよ。問題~\ref{ex:08:4}の見積もりはどこで正しいか。変化が速いとき、$g^*$が減らなくなるのはなぜか。
\end{exercise}

\begin{exercise}[\lvlB]
\label{ex:08:6}
\emph{設計：選択ネットワークへの割り込み。}問題~\ref{ex:08:2}のネットワークは$\beta=2$、$\tau=20\ms$、$g=1$で、入力が今は$I_1=0.9$、$I_2=1.0$なのに$r_1=0.9$、$r_2=0$を保っている。\nt{NORA}は時間$T_p$の間$g$を$g_{\text{lo}}$に下げる。$g_{\text{lo}}=0$のときの最小の$T_p$を解析的に、$g_{\text{lo}}=0.25$のときをモデルで求めよ。
\end{exercise}

\begin{exercise}[\lvlC]
\label{ex:08:7}
\emph{調整が割に合うとき。}オラクルは穏やかな時期か変化の多い時期かを知っていて、それぞれに自分の$g$を設定する。(a)~本章の世界（各時期$1000$試行、$h=0.001$と$0.05$）で、最良の一定$g$に対するその利得をモデルで求めよ。(b)~利得がより大きい世界を構成し、驚きによる調整がそのどれだけを取り戻すかを求めよ。
\end{exercise}
